\chapter{Inguinal Hernia}

\section*{Definition and Overview}
An inguinal hernia is a bulge in the groin caused by part of the abdominal contents, usually a loop of intestine or fatty tissue, pushing through a weakness in the lower abdominal wall. It is the most common type of hernia and occurs much more often in men, in whom the passage taken by the testicle in development leaves a natural point of weakness. Some inguinal hernias cause no symptoms, while others produce a visible bulge and discomfort, especially with coughing, straining, or lifting. Small or symptom-free hernias may be safely monitored, but hernias that are painful, enlarging, or strangulating require repair. Surgery is a common, low-risk procedure with excellent outcomes.

\section*{Epidemiology}
Inguinal hernias are very common, affecting around one in four men and one in twenty women at some stage of life. The lifetime risk is far higher in men, and incidence increases with age. Risk factors include a family history of hernia, chronic coughing from smoking or lung disease, constipation with straining, heavy lifting, obesity, and conditions that raise abdominal pressure, such as an enlarged prostate or pregnancy. Repair of inguinal hernia is one of the most frequently performed operations in Australia and worldwide, with many thousands performed each year.

\section*{Aetiology and Pathophysiology}
In men, the inguinal canal is the passage through which the spermatic cord travels to the scrotum; in women, it carries the round ligament. Weakness at this site allows abdominal contents to protrude, forming a hernia. Direct hernias push through a weak area of the abdominal wall itself, usually in older men, while indirect hernias follow the inguinal canal and are the common type in younger people, and in babies, where the canal has not closed after the testicle descends. Raised intra-abdominal pressure from coughing, straining, lifting, or obesity forces contents through the defect. Complications arise when the contents become trapped (incarcerated) and when their blood supply is cut off (strangulated), which is a surgical emergency.

\section*{Clinical Features}
\begin{itemize}
    \item A bulge in the groin that may appear with standing, coughing, or straining
    \item A feeling of heaviness, dragging, or discomfort in the groin
    \item The bulge often reduces (goes back in) when lying down
    \item Pain that increases with lifting, coughing, or prolonged standing
    \item In men, the bulge may extend into the scrotum
    \item A sudden, painful, and non-reducible bulge, with nausea and vomiting, may indicate strangulation
    \item Many hernias cause no symptoms and are found on examination
\end{itemize}

\section*{Differential Diagnosis}
Not every groin bulge is an inguinal hernia.
\begin{itemize}
    \item \textbf{Femoral hernia:} a lower groin bulge, more common in women and more prone to strangulation
    \item \textbf{Hydrocele:} a fluid collection in the scrotum that transilluminates
    \item \textbf{Undescended testicle or testicular swelling:} in children and men
    \item \textbf{Enlarged lymph nodes:} from infection or malignancy
    \item \textbf{Spermatic cord cysts and varicoceles:} other scrotal or groin swellings
    \item \textbf{Muscle strain:} groin pain without a bulge
    \item \textbf{Hip joint problems:} pain radiating into the groin
\end{itemize}

\section*{When to Seek Medical Care}
Anyone with a new or persistent groin bulge should have it assessed, and men and women with risk factors should be examined if symptoms develop. Repair is recommended for hernias that are painful, enlarging, or affecting work and daily life. Even symptom-free hernias need regular review, because they can enlarge or become complicated.

\textbf{Call 000 (or local emergency services) for:}
\begin{itemize}
    \item A painful, tense bulge that will not go back in, especially with vomiting, severe pain, or abdominal distension, which may indicate strangulation
    \item Sudden severe groin or abdominal pain
    \item Fever, rapid heart rate, or signs of serious illness with a hernia
    \item A tender bulge in a baby or child with vomiting or distress
\end{itemize}

\section*{Diagnosis and Investigations}
\begin{itemize}
    \item \textbf{History and examination:} the bulge is usually confirmed by examination, with the person standing and coughing
    \item \textbf{Assessment of reducibility:} whether the bulge returns with gentle pressure or lying down
    \item \textbf{Ultrasound:} helpful when the diagnosis is uncertain or in overweight people
    \item \textbf{CT scan:} for complex, large, or recurrent hernias, or when complications are suspected
    \item \textbf{Urgent assessment:} for signs of incarceration or strangulation
    \item \textbf{Preoperative assessment:} fitness for anaesthesia before planned repair
\end{itemize}

\section*{Management}
\begin{itemize}
    \item \textbf{Watchful waiting:} symptom-free or mildly symptomatic hernias can be safely monitored, with repair if symptoms develop
    \item \textbf{Open hernia repair:} the hernia is returned and the defect closed, usually with a synthetic mesh to reinforce the wall
    \item \textbf{Laparoscopic repair:} a minimally invasive technique through small incisions, with faster recovery in suitable cases
    \item \textbf{Emergency surgery:} for incarcerated or strangulated hernias, with the bowel assessed for viability
    \item \textbf{Pain relief and symptom management:} while awaiting or deciding about surgery
    \item \textbf{Addressing risk factors:} treating chronic cough, constipation, and lifting practices
    \item \textbf{Aftercare:} graded return to activity with lifting restrictions as advised
\end{itemize}

\section*{Complications}
\begin{itemize}
    \item Incarceration: the hernia contents become trapped
    \item Strangulation: the blood supply is cut off, with bowel death and perforation, a life-threatening emergency
    \item Bowel obstruction
    \item Complications of surgery: bleeding, infection, chronic pain, and recurrence
    \item Injury to the spermatic cord structures, affecting fertility in rare cases
    \item Mesh complications, including infection and chronic discomfort
\end{itemize}

\section*{Prognosis and Outcomes}
Modern hernia repair is highly successful, with recurrence rates of around 1 to 5 per cent depending on technique and risk factors. Most people return to normal activities within a few weeks, and laparoscopic repair offers a somewhat faster recovery. For symptom-free hernias, watchful waiting is safe, and around a third will eventually need surgery as they enlarge or become symptomatic. The serious complications of strangulation are largely preventable by timely repair.

\section*{Prevention and Self-Care}
\begin{itemize}
    \item Maintain a healthy weight
    \item Avoid chronic straining: treat constipation and lift with the legs rather than the back
    \item Stop smoking to reduce chronic cough
    \item Seek treatment for an enlarged prostate or chronic lung disease
    \item Have any new groin bulge assessed
    \item Consider elective repair before symptoms become severe
    \item After surgery, follow lifting and activity advice and seek review for any new bulge or pain
\end{itemize}

\section*{Sources and Further Reading}
The following reputable sources provide further information for healthcare students and patients seeking reliable, current advice about inguinal hernia.
\begin{itemize}
    \item Healthdirect Australia. \textit{Inguinal hernia.} https://www.healthdirect.gov.au/
    \item Royal Australasian College of Surgeons. \textit{Hernia repair information.} https://www.surgeons.org/
    \item Better Health Channel. \textit{Hernias.} https://www.betterhealth.vic.gov.au/
    \item NHS. \textit{Inguinal hernia repair.} https://www.nhs.uk/conditions/inguinal-hernia-repair/
    \item Mayo Clinic. \textit{Inguinal hernia.} https://www.mayoclinic.org/
    \item MSD Manual. \textit{Hernias of the abdominal wall.} https://www.msdmanuals.com/
\end{itemize}
